% GENERATED FILE. Edit canonical lessons, then run npm run generate:books.
% canonical-source-hash: fnv1a64:9062e967b01eaf2a
% canonical-lessons: MW-C281-chak, MW-C281-ghumavno, MW-C281-thapno, MW-C281-korno, MW-C281-bhatti

\chapter{Clay, Wood and Mud}
\label{ch:mw-clay-wood-and-mud}
\hlchaptermodality{281}

\begin{chapteropening}
\textbf{By the end of this chapter:} \emph{I can name the potter's and the carver's work.}

Complete the last lesson of chapter 281: I can name the potter's and the carver's work.
\end{chapteropening}

\section[\texorpdfstring{\mw{चाक} (chāk)}{चाक (chāk)}]{\mw{चाक} (chāk) — a potter's wheel}

\label{lesson:MW-C281-chak}



\begin{quote}
Before the new one: say the Marwadi for a target, then the Marwadi for to take aim.
\end{quote}

\subsection*{You'll want to know: \mw{चाक}}
\textbf{\mw{चाक}} — \emph{chāk} — \textquotedblleft{}a potter's wheel\textquotedblright{}.

\begin{quote}
\mw{चाक} \mw{घणो} \mw{भारी} \mw{है।}
\end{quote}

\begin{quote}
\emph{chāk ghaṇo bhārī hai.}
\end{quote}

\begin{quote}
The potter's wheel is very heavy.
\end{quote}

A masculine noun.

\begin{cousinweb}
\textbf{\mw{चाक}} (\emph{chāk}) is inherited from Sanskrit \emph{cakra}, \textquotedblleft{}a wheel; a potter's wheel\textquotedblright{}; CDIAL gives no Marwari form in this entry, but lists Hindi \emph{cāk}, \textquotedblleft{}any kind of wheel\textquotedblright{}, and Gujarati \emph{cāk}, \textquotedblleft{}wheel, potter's wheel\textquotedblright{} (CDIAL 4538).
\end{cousinweb}

\begin{grammarlens}[title={What you've built}]
One more word for this chapter.
\end{grammarlens}

\subsection*{Your turn}
\begin{itemize}
\raggedright
  \item \emph{Say it:} \emph{chāk}
  \item \emph{Say it:} \emph{chāk}, once more
  \item \emph{Say it:} \emph{chāk ghaṇo bhārī hai.} — The potter's wheel is very heavy.
  \item \emph{Recall:} say the Marwadi for a target, then the Marwadi for to take aim, then say \emph{chāk} again
\end{itemize}

\begin{tcolorbox}[breakable,colback=teal!4,colframe=teal!35!black,title={Before you move on}]
What does \mw{चाक} mean? (\textquotedblleft{}A potter's wheel\textquotedblright{}.) Say it once more.
\end{tcolorbox}
\section[\texorpdfstring{\mw{घुमावणो} (ghumāvṇo)}{घुमावणो (ghumāvṇo)}]{\mw{घुमावणो} (ghumāvṇo) — to turn (a wheel); to spin}

\label{lesson:MW-C281-ghumavno}



\begin{quote}
Before the new one: say the Marwadi for to take aim, then the Marwadi for a potter's wheel.
\end{quote}

\subsection*{You'll want to know: \mw{घुमावणो}}
\textbf{\mw{घुमावणो}} — \emph{ghumāvṇo} — \textquotedblleft{}to turn (a wheel); to spin\textquotedblright{}.

\begin{quote}
\mw{काको} \mw{चाक} \mw{घुमावै} \mw{है।}
\end{quote}

\begin{quote}
\emph{kāko chāk ghumāvai hai.}
\end{quote}

\begin{quote}
My uncle turns the potter's wheel.
\end{quote}

Past (I …): \mw{म्हैं} \mw{चाक} \mw{घुमायो}, \emph{mhaĩ chāk ghumāyo}.

\begin{cousinweb}
\textbf{\mw{घुमावणो}} (\emph{ghumāvṇo}) is the causative of \textbf{\mw{घूमणो}}, to go round: it is built on that verb and means \textquotedblleft{}to make go round\textquotedblright{}. CDIAL gives no Marwari form in this entry, but lists Hindi \emph{ghumnā}, \textquotedblleft{}to revolve\textquotedblright{}, under a reconstructed \emph{*ghummati} (CDIAL 4485).
\end{cousinweb}

\begin{grammarlens}[title={What you've built}]
One more word for this chapter.
\end{grammarlens}

\subsection*{Your turn}
\begin{itemize}
\raggedright
  \item \emph{Say it:} \emph{ghumāvṇo}
  \item \emph{Say it:} \emph{ghumāvṇo}, once more
  \item \emph{Say it:} \emph{kāko chāk ghumāvai hai.} — My uncle turns the potter's wheel.
  \item \emph{Recall:} say the Marwadi for to take aim, then the Marwadi for a potter's wheel, then say \emph{ghumāvṇo} again
\end{itemize}

\begin{tcolorbox}[breakable,colback=teal!4,colframe=teal!35!black,title={Before you move on}]
What does \mw{घुमावणो} mean? (\textquotedblleft{}To turn (a wheel); to spin\textquotedblright{}.) Say it once more.
\end{tcolorbox}
\section[\texorpdfstring{\mw{थापणो} (thāpṇo)}{थापणो (thāpṇo)}]{\mw{थापणो} (thāpṇo) — to pat into shape}

\label{lesson:MW-C281-thapno}



\begin{quote}
Before the new one: say the Marwadi for a potter's wheel, then the Marwadi for to spin a wheel.
\end{quote}

\subsection*{You'll want to know: \mw{थापणो}}
\textbf{\mw{थापणो}} — \emph{thāpṇo} — \textquotedblleft{}to pat into shape\textquotedblright{}.

\begin{quote}
\mw{मां} \mw{छाणो} \mw{थापै} \mw{है।}
\end{quote}

\begin{quote}
\emph{mā̃ chhāṇo thāpai hai.}
\end{quote}

\begin{quote}
Mother pats out a dung cake.
\end{quote}

Past (I …): \mw{म्हैं} \mw{छाणो} \mw{थाप्यो}, \emph{mhaĩ chhāṇo thāpyo}.

\begin{cousinweb}
\textbf{\mw{थापणो}} (\emph{thāpṇo}) is inherited from a reconstructed \emph{*thapp}, \textquotedblleft{}slap, pat\textquotedblright{}; CDIAL gives no Marwari form in this entry, but lists Hindi \emph{thāpnā}, \textquotedblleft{}to pat\textquotedblright{}, and Gujarati \emph{thāpvũ}, \textquotedblleft{}to flatten by patting\textquotedblright{} (CDIAL 6091).
\end{cousinweb}

\begin{grammarlens}[title={What you've built}]
One more word for this chapter.
\end{grammarlens}

\subsection*{Your turn}
\begin{itemize}
\raggedright
  \item \emph{Say it:} \emph{thāpṇo}
  \item \emph{Say it:} \emph{thāpṇo}, once more
  \item \emph{Say it:} \emph{mā̃ chhāṇo thāpai hai.} — Mother pats out a dung cake.
  \item \emph{Recall:} say the Marwadi for a potter's wheel, then the Marwadi for to spin a wheel, then say \emph{thāpṇo} again
\end{itemize}

\begin{tcolorbox}[breakable,colback=teal!4,colframe=teal!35!black,title={Before you move on}]
What does \mw{थापणो} mean? (\textquotedblleft{}To pat into shape\textquotedblright{}.) Say it once more.
\end{tcolorbox}
\section[\texorpdfstring{\mw{कोरणो} (korṇo)}{कोरणो (korṇo)}]{\mw{कोरणो} (korṇo) — to carve}

\label{lesson:MW-C281-korno}



\begin{quote}
Before the new one: say the Marwadi for to spin a wheel, then the Marwadi for to pat into shape.
\end{quote}

\subsection*{You'll want to know: \mw{कोरणो}}
\textbf{\mw{कोरणो}} — \emph{korṇo} — \textquotedblleft{}to carve\textquotedblright{}.

\begin{quote}
\mw{काको} \mw{लकड़ी} \mw{कोरै} \mw{है।}
\end{quote}

\begin{quote}
\emph{kāko lakṛī korai hai.}
\end{quote}

\begin{quote}
My uncle carves wood.
\end{quote}

Past (I …): \mw{म्हैं} \mw{लकड़ी} \mw{कोरी}, \emph{mhaĩ lakṛī korī}.

\begin{cousinweb}
\textbf{\mw{कोरणो}} (\emph{korṇo}) is inherited from a reconstructed \emph{*kōrati}, \textquotedblleft{}scrapes\textquotedblright{}; CDIAL gives no Marwari form in this entry, but lists Hindi \emph{kornā}, \textquotedblleft{}to scrape or scoop out\textquotedblright{}, and Gujarati \emph{korvũ}, \textquotedblleft{}to scoop, carve\textquotedblright{} (CDIAL 3530).
\end{cousinweb}

\begin{grammarlens}[title={What you've built}]
One more word for this chapter.
\end{grammarlens}

\subsection*{Your turn}
\begin{itemize}
\raggedright
  \item \emph{Say it:} \emph{korṇo}
  \item \emph{Say it:} \emph{korṇo}, once more
  \item \emph{Say it:} \emph{kāko lakṛī korai hai.} — My uncle carves wood.
  \item \emph{Recall:} say the Marwadi for to spin a wheel, then the Marwadi for to pat into shape, then say \emph{korṇo} again
\end{itemize}

\begin{tcolorbox}[breakable,colback=teal!4,colframe=teal!35!black,title={Before you move on}]
What does \mw{कोरणो} mean? (\textquotedblleft{}To carve\textquotedblright{}.) Say it once more.
\end{tcolorbox}
\section[\texorpdfstring{\mw{भट्टी} (bhaṭṭī)}{भट्टी (bhaṭṭī)}]{\mw{भट्टी} (bhaṭṭī) — a kiln}

\label{lesson:MW-C281-bhatti}



\begin{quote}
Before the new one: say the Marwadi for to pat into shape, then the Marwadi for to carve.
\end{quote}

\subsection*{You'll want to know: \mw{भट्टी}}
\textbf{\mw{भट्टी}} — \emph{bhaṭṭī} — \textquotedblleft{}a kiln\textquotedblright{}.

\begin{quote}
\mw{भट्टी} \mw{घणी} \mw{गरम} \mw{है।}
\end{quote}

\begin{quote}
\emph{bhaṭṭī ghaṇī garam hai.}
\end{quote}

\begin{quote}
The kiln is very hot.
\end{quote}

A feminine noun.

\begin{cousinweb}
\textbf{\mw{भट्टी}} (\emph{bhaṭṭī}) is inherited from Sanskrit \emph{bhrāṣṭra}, \textquotedblleft{}a frying pan\textquotedblright{}; CDIAL gives no Marwari form in this entry, but lists Marathi \emph{bhaṭṭī}, \textquotedblleft{}forge\textquotedblright{}, Sindhi \emph{baṭhī}, \textquotedblleft{}distilling furnace\textquotedblright{}, and Hindi \emph{bhaṭṭhā}, \textquotedblleft{}kiln\textquotedblright{} (CDIAL 9656).
\end{cousinweb}

\begin{grammarlens}[title={What you've built}]
One more word for this chapter.
\end{grammarlens}

\subsection*{Your turn}
\begin{itemize}
\raggedright
  \item \emph{Say it:} \emph{bhaṭṭī}
  \item \emph{Say it:} \emph{bhaṭṭī}, once more
  \item \emph{Say it:} \emph{bhaṭṭī ghaṇī garam hai.} — The kiln is very hot.
  \item \emph{Recall:} say the Marwadi for to pat into shape, then the Marwadi for to carve, then say \emph{bhaṭṭī} again
\end{itemize}

\begin{tcolorbox}[breakable,colback=teal!4,colframe=teal!35!black,title={Before you move on}]
What does \mw{भट्टी} mean? (\textquotedblleft{}A kiln\textquotedblright{}.) Say it once more.
\end{tcolorbox}
